% ============================================================================
\chapter{Arquitetura do Sistema}
% ============================================================================

Este capítulo apresenta a arquitetura proposta para o EduHub a partir das
restrições, premissas e requisitos não funcionais definidos no Capítulo 1. A
arquitetura é modelada segundo o
padrão C4, com diagramas de contexto (Nível 1), de contêineres (Nível 2) e de
componentes (Nível 3). As decisões descritas derivam diretamente dessas
restrições e requisitos.

\section{Contexto e fronteiras do sistema}

O EduHub é um sistema autônomo e independente,
que não é componente de um sistema maior e não depende de sistemas acadêmicos
legados nem de gateways de pagamento durante a operação. Sua única integração
externa é com um serviço de e-mail, usado para o envio de notificações
(RF-36). Os dados primários do sistema são inseridos e mantidos pela própria
interface administrativa e, após a implantação, o EduHub passa a ser a fonte
primária dos dados acadêmicos da instituição.

A Figura~\ref{fig:c4-contexto} apresenta o diagrama de contexto (C4, Nível
1). A fronteira do sistema compreende o próprio EduHub, os três perfis de
usuário (Administrativo, Professor e Aluno), descritos na
Tabela~\ref{tab:perfis-usuarios}, e o serviço de e-mail externo. A única
outra dependência externa é operacional, a infraestrutura de servidor de
produção fornecida pela equipe de TI da instituição (D-1, ver
Seção~\ref{sec:restricoes-arquitetura}), e não uma integração de dados.

\begin{figure}[ht]
\centering
\includegraphics[width=\textwidth]{c4-contexto.png}
\caption{Diagrama de contexto do EduHub (C4 -- Nível 1). Fonte: elaboração própria da equipe.}
\label{fig:c4-contexto}
\end{figure}

\section{Estilo ou padrão arquitetural}

O EduHub adota uma organização modular por domínio funcional, derivada do
RNF-07. O sistema é estruturado em módulos que correspondem às áreas de
funcionalidade identificadas no Capítulo 1: Acesso e Autenticação, Cadastros
e Gestão da Estrutura, Gestão Acadêmica, Informação e Consulta do Aluno,
Comunicação e Notificações, e Relatórios e Auditoria. Cada módulo encapsula
suas próprias regras de negócio e dados, o que reduz o acoplamento entre
áreas que evoluem em ritmos diferentes. Por exemplo, o módulo de Comunicação
pode receber novos recursos sem alterar o módulo de Gestão Acadêmica.

Essa organização modular sustenta o objetivo de qualidade de
manutenibilidade, apontado na Seção~\ref{sec:quality-goals} como prioridade
alta, uma vez que módulos com fronteiras claras favorecem alta coesão e baixo
acoplamento.

\section{Contêineres e componentes}

A arquitetura do EduHub foi modelada seguindo o padrão C4. Os diagramas de
contêineres (Nível 2) e de componentes (Nível 3) detalham a estrutura interna
do sistema a partir da visão de contexto.

No Nível 2 (contêineres), ilustrado na Figura~\ref{fig:c4-container}, o EduHub
é decomposto em três contêineres:
\begin{itemize}[nosep]
    \item \textbf{Frontend}: aplicação web responsiva acessada pelos usuários
    via navegador, que concentra as telas de login, cadastro, grade horária e
    demais interfaces.
    \item \textbf{Backend}: responsável pela lógica de negócio e pela API que
    atende às requisições do Frontend por meio de chamadas JSON/HTTP.
    \item \textbf{Banco de Dados}: armazena os dados de alunos, disciplinas e
    demais informações acadêmicas.
\end{itemize}
O Backend lê e escreve no Banco de Dados e solicita ao serviço de e-mail
externo o envio das notificações.

\begin{figure}[ht]
\centering
\includegraphics[width=\textwidth]{c4-container.png}
\caption{Diagrama de contêineres do EduHub (C4 -- Nível 2). Fonte: elaboração própria da equipe.}
\label{fig:c4-container}
\end{figure}

Cada contêiner do Backend concentra os módulos funcionais do sistema. A
Tabela~\ref{tab:modulos} consolida esses módulos e suas responsabilidades,
servindo de ponte entre os requisitos e os componentes detalhados a seguir.

\begin{table}[H]
\centering
\caption{Módulos funcionais do EduHub.}
\label{tab:modulos}
\begin{tabularx}{\textwidth}{p{0.28\textwidth} X}
\toprule
\textbf{Módulo} & \textbf{Responsabilidade principal} \\
\midrule
Acesso e Autenticação & Autenticação de usuários por credencial, recuperação
de senha, gestão de cargos, hierarquia e permissões (RF-01 a RF-13). \\
Cadastros e Gestão da Estrutura & Cadastro de dados da escola, turmas,
disciplinas, horários, calendário acadêmico e vínculo professor-disciplina
(RF-14 a RF-20). \\
Gestão Acadêmica & Lançamento e edição de notas, registro de frequência,
cálculo automático de faltas e travamento de períodos letivos (RF-21 a
RF-27). \\
Informação e Consulta do Aluno & Consulta a boletim, frequência, grade de
horários, calendário e solicitação de documentos pelo aluno (RF-28 a
RF-34). \\
Comunicação e Notificações & Envio de comunicados pela administração,
notificações por e-mail, caixa de entrada e mensagens do aluno com roteamento
por tag de assunto (RF-35 a RF-41). \\
Relatórios e Auditoria & Geração e exportação de relatórios acadêmicos e
consulta ao log de alterações e auditoria (RF-42 a RF-45). \\
\bottomrule
\end{tabularx}
\end{table}

No Nível 3 (componentes), ilustrado na Figura~\ref{fig:c4-componentes}, o
contêiner Backend é detalhado em oito componentes que correspondem, em
conjunto, aos módulos funcionais da Tabela~\ref{tab:modulos}:
\begin{itemize}[nosep]
    \item \textbf{Autenticação e Acesso}: login, 2FA e permissões.
    \item \textbf{Gestão de Usuários}: cadastro de alunos, professores e
    administradores.
    \item \textbf{Estrutura Escolar}: cadastro base de escola, turmas e
    disciplinas.
    \item \textbf{Autoatendimento do Aluno}: consulta de boletim, grade
    horária e calendário.
    \item \textbf{Registro Acadêmico}: lançamento de notas e frequência.
    \item \textbf{Mensageria e Comunicação}: caixa de entrada do aluno e
    envio de avisos.
    \item \textbf{Adaptador de E-mail}: formata e dispara os e-mails.
    \item \textbf{Relatório e Auditoria}: geração de relatórios e gravação de
    logs.
\end{itemize}
Os componentes gravam seus dados no Banco de Dados. O Adaptador de E-mail
isola a comunicação com o serviço de e-mail externo, recebendo o e-mail não
formatado da Mensageria, formatando-o e disparando-o para o serviço externo.

\begin{figure}[ht]
\centering
\includegraphics[width=\textwidth]{c4-componentes.png}
\caption{Diagrama de componentes do Backend do EduHub (C4 -- Nível 3). Fonte: elaboração própria da equipe.}
\label{fig:c4-componentes}
\end{figure}


\section{Decisões arquiteturais relevantes}

A tabela a seguir registra as decisões arquiteturais que decorrem
diretamente das restrições e requisitos não funcionais especificados no
Capítulo 1.

\begin{table}[ht]
\centering
\caption{Decisões arquiteturais do EduHub.}
\begin{tabularx}{\textwidth}{p{0.24\textwidth} X X}
\toprule
\textbf{Decisão} & \textbf{Motivação} & \textbf{Consequências} \\
\midrule
Aplicação web responsiva, sem aplicativo móvel nativo &
R-1 exige acesso via navegador sem app nativo dedicado nesta versão &
Reduz o esforço de desenvolvimento e manutenção (uma única base de código
para todos os dispositivos), mas a experiência em telas pequenas depende
inteiramente do design responsivo. \\
Arquitetura modular por domínio funcional &
RNF-07 exige que novos recursos sejam adicionados sem impactar módulos
existentes &
Facilita manutenção e evolução incremental (alinhado ao processo Scrumban do
projeto), mas exige disciplina para manter baixo acoplamento entre módulos
desde as primeiras iterações. \\
Conformidade com a LGPD e criptografia de dados sensíveis &
R-3 exige atendimento à legislação de proteção de dados e armazenamento
seguro de credenciais e dados sensíveis. O RNF-03 detalha a criptografia em
repouso de dados cadastrais, senhas e e-mails &
Aumenta a complexidade do módulo de Acesso e Autenticação e da camada de
persistência, mas é indispensável para tratar dados pessoais de alunos e
professores. \\
Suporte opcional a autenticação de dois fatores (2FA) &
RNF-04 exige suporte opcional a 2FA para os perfis administrativo e docente &
Reforça a segurança de acesso para os perfis com maior privilégio, sem
tornar o login obrigatoriamente mais complexo para o perfil aluno. \\
Acessibilidade conforme WCAG 2.1 (Nível AA) &
R-2 e RNF-10 exigem que a interface siga as diretrizes WCAG 2.1 AA &
Restringe escolhas de componentes de interface e exige testes de
acessibilidade, mas garante inclusão de usuários com deficiência entre os
três perfis. \\
Sistema autônomo, com única integração externa (serviço de e-mail) &
A equipe definiu o EduHub, já na especificação de requisitos (item 1.2.1),
como sistema independente e fonte primária dos próprios dados. A única
dependência externa em operação é o serviço de envio de e-mail para
notificações (RF-36) &
Simplifica a arquitetura (sem integrações com sistemas acadêmicos legados ou
gateways de pagamento) e isola a dependência externa num único componente
(Adaptador de E-mail). Dados de sistemas legados, se existirem, precisam ser
inseridos manualmente pelo perfil administrativo. \\
\bottomrule
\end{tabularx}
\end{table}

\section{Restrições de arquitetura}
\label{sec:restricoes-arquitetura}

As restrições, premissas e dependências a seguir foram definidas pela equipe
já na etapa de especificação de requisitos (itens 1.2.4 a 1.2.6) e limitam
diretamente as decisões de arquitetura e implementação do EduHub.

\begin{table}[ht]
\centering
\caption{Restrições, premissas e dependências do projeto.}
\begin{tabularx}{\textwidth}{p{0.10\textwidth} X}
\toprule
\textbf{ID} & \textbf{Descrição} \\
\midrule
R-1 & O sistema deve ser desenvolvido como uma aplicação web responsiva
(acessível via navegador), sem a necessidade de um aplicativo móvel nativo
dedicado nesta versão. \\
R-2 & A interface do usuário deve aderir aos padrões de acessibilidade WCAG
2.1 (Nível AA) para garantir a inclusão de usuários com deficiência. \\
R-3 & O sistema deve atender à legislação de proteção de dados (LGPD) no
tratamento de informações pessoais de alunos e professores, e deve garantir
que todas as credenciais de acesso e dados sensíveis sejam armazenados
utilizando métodos criptográficos seguros. \\
\midrule
P-1 & Conectividade: presume-se que professores e administrativo terão
acesso estável à internet com banda larga na instituição de ensino para a
inserção de dados. \\
P-2 & Navegadores suportados: presume-se uso restrito a Google Chrome e
derivados Chromium, Mozilla Firefox e Safari (versões mais recentes); outros
navegadores estão fora do escopo de teste e suporte. \\
P-3 & Proficiência do usuário: presume-se que os usuários administrativos
possuem proficiência básica no uso de computadores e ferramentas de
escritório. \\
P-4 & Processo de validação de notas: presume-se que os professores são
responsáveis pela validação final e travamento das notas, cabendo à
administração apenas monitorar esse processo. \\
\midrule
D-1 & Ambiente de produção: o projeto depende da alocação e configuração,
pela equipe de TI da instituição, de um servidor de produção com as
especificações mínimas de hardware e sistema operacional. \\
\bottomrule
\end{tabularx}
\end{table}

% ============================================================================
\chapter{Projeto de Componentes}
% ============================================================================

Este capítulo detalha, a partir dos módulos funcionais e dos componentes
apresentados no Capítulo 2, os principais componentes de software do EduHub e
as regras de negócio associadas a cada um, com base nos requisitos funcionais
definidos no Capítulo 1. Para os componentes de Autenticação e Acesso e de
Registro Acadêmico, dois dos mais centrais do sistema, apresenta-se também o
projeto de classes correspondente na Seção 3.2.

\section{Visão geral dos componentes}

A Tabela~\ref{tab:componentes} lista os principais componentes do sistema,
que correspondem aos componentes do diagrama C4 de Nível 3
(Figura~\ref{fig:c4-componentes}) e aos requisitos funcionais que cada um
implementa.

\begin{table}[ht]
\centering
\caption{Principais componentes do EduHub.}
\label{tab:componentes}
\begin{tabularx}{\textwidth}{p{0.30\textwidth} X}
\toprule
\textbf{Componente} & \textbf{Responsabilidade principal} \\
\midrule
Autenticação e Acesso & Login por credencial, recuperação de senha, registro
de histórico de login, autenticação de dois fatores e gestão de cargos,
hierarquia e permissões (RF-01, RF-02, RF-05 a RF-10). \\
Gestão de Usuários & Cadastro, edição e inativação de alunos, professores e
administradores, incluindo dados de perfil (RF-03, RF-04, RF-11 a RF-13). \\
Estrutura Escolar & Cadastro base da escola, turmas, disciplinas, horários,
calendário acadêmico, vínculo professor-disciplina e prevenção de conflitos
de horário (RF-14 a RF-20). \\
Registro Acadêmico & Lançamento e edição de notas, registro de frequência,
cálculo automático de faltas e travamento de lançamentos por período letivo
(RF-21 a RF-27). \\
Autoatendimento do Aluno & Consulta de boletim, frequência detalhada, grade
de horários, calendário e solicitação de documentos (RF-28 a RF-34). \\
Mensageria e Comunicação & Envio de comunicados pela administração, caixa de
entrada do aluno e roteamento de mensagens por tag de assunto (RF-35 a
RF-41). \\
Adaptador de E-mail & Formatação e disparo das notificações por e-mail para o
serviço externo (RF-36). \\
Relatório e Auditoria & Geração e exportação de relatórios acadêmicos e
registro e consulta dos logs de alteração (RF-42 a RF-45). \\
\bottomrule
\end{tabularx}
\end{table}

\section{Detalhamento dos componentes principais}
\label{sec:detalhamento}

Dentre os componentes listados, dois são destacados por concentrarem regras
de negócio mais complexas e por serem centrais aos objetivos de segurança e
qualidade acadêmica do sistema.

\textbf{Autenticação e Acesso:} Este componente concentra a autenticação dos
usuários e o controle de cargos e permissões, garantindo que todo usuário do
sistema opere sob um cargo bem definido. As regras de negócio associadas são:
\begin{itemize}[nosep]
    \item qualquer usuário se autentica por e-mail e senha (RF-01) e pode
    solicitar a recuperação de senha por e-mail (RF-02), com registro do
    histórico de login para auditoria (RF-05).
    \item o sistema oferece suporte opcional a autenticação de dois fatores
    para os perfis administrativo e docente (RNF-04).
    \item todo novo usuário deve ser obrigatoriamente associado a um cargo no
    momento da criação (RF-06).
    \item apenas usuários com permissão administrativa podem designar ou
    alterar o cargo de outro usuário (RF-07), e a alteração atualiza
    imediatamente as permissões de acesso do usuário-alvo, conforme o CU-01.
    \item o usuário administrativo pode criar cargos personalizados e definir
    seu nível de hierarquia (RF-08).
    \item o sistema mantém cargos padrões que não podem ser excluídos, mas
    cujo nome de exibição pode ser alterado (RF-09).
    \item permissões específicas podem ser definidas por cargo, seja ele
    padrão ou personalizado (RF-10).
\end{itemize}

O projeto de classes correspondente ao componente de Autenticação e Acesso é
apresentado na Figura~\ref{fig:classes-autenticacao}.

\begin{figure}[H]
\centering
\includegraphics[width=\textwidth]{classes-autenticacao-acesso.png}
\caption{Diagrama de classes do componente de Autenticação e Acesso do EduHub. Fonte: elaboração própria da equipe.}
\label{fig:classes-autenticacao}
\end{figure}

O diagrama organiza as entidades
centrais do domínio de acesso, como \texttt{Usuario}, \texttt{Cargo},
\texttt{Permissao} e \texttt{HistoricoLogin}, em que cada usuário é associado a
um único cargo e cada cargo agrupa as permissões que definem o que seus
usuários podem fazer no sistema (RF-06, RF-10). A autenticação de dois fatores
(RNF-04) é modelada pela interface \texttt{Metodo2FA}, com as operações
\texttt{gerarCodigo()} e \texttt{validarCodigo()}, implementada pelas classes
concretas \texttt{TwoFactorEmail} e \texttt{TwoFactorApp}, o que permite
acrescentar novos meios de verificação sem alterar o restante do componente.
As regras de negócio são expostas pela interface \texttt{AutenticacaoService},
implementada por \texttt{AutenticacaoServiceImpl}, com operações como
\texttt{autenticar()}, \texttt{recuperarSenha()}, \texttt{atribuirCargo()} e
\texttt{verificarPermissao()}. A persistência é isolada pela interface
\texttt{Autenticacao\allowbreak Repository}, responsável por operações como
\texttt{salvarUsuario()}, \texttt{buscarUsuarioPorEmail()},
\texttt{salvarCargo()} e \texttt{registrarHistoricoLogin()}. Assim como no
Registro Acadêmico, essa separação entre serviço e repositório mantém as
regras de negócio independentes da tecnologia de persistência, em linha com o
objetivo de manutenibilidade (RNF-07).

\textbf{Registro Acadêmico:} Este componente concentra as
regras de negócio do processo de avaliação e controle de presença, com
forte dependência do conceito de período letivo travado/destravado
(RF-26, premissa P-4):
\begin{itemize}[nosep]
    \item o professor tem acesso exclusivo à lista de alunos e à estrutura
    de notas apenas das turmas e disciplinas às quais está vinculado
    (RF-21).
    \item o professor pode inserir notas parciais e finais por disciplina e
    por período letivo (RF-22) e editá-las (RF-23), desde que o período não
    tenha sido travado pela administração (RF-23, RF-26).
    \item o professor registra presença, falta ou justificativa de ausência
    por aluno, data e aula/período (RF-24).
    \item o sistema calcula automaticamente o total de faltas e o
    percentual de frequência por disciplina (RF-25).
    \item o usuário administrativo pode travar o lançamento de notas e
    frequência por período letivo a qualquer momento, impedindo qualquer
    alteração posterior pelos professores (RF-26). Conforme CU-03, se houver
    turmas ou disciplinas do período com notas finais ou frequência
    pendentes, o sistema apenas sinaliza essas pendências ao administrativo,
    sem impedir o travamento. Uma correção após o período travado exige que
    o professor solicite o destravamento à administração, que reabre o
    período pela função ``Desbloquear Período'', permite a correção e trava
    o período novamente.
    \item o professor pode visualizar, a partir da lista de chamada, os
    dados de contato do aluno cadastrados pelo administrativo (RF-27).
\end{itemize}

O projeto de classes correspondente ao componente de Registro Acadêmico é
apresentado na Figura~\ref{fig:classes-registro}.

\begin{figure}[ht]
\centering
\includegraphics[width=\textwidth]{classes-registro-academico.png}
\caption{Diagrama de classes do componente de Registro Acadêmico do EduHub. Fonte: elaboração própria da equipe.}
\label{fig:classes-registro}
\end{figure}

O diagrama organiza as entidades centrais
do domínio acadêmico, como \texttt{Nota}, \texttt{Aluno}, \texttt{Turma},
\texttt{Disciplina} e \texttt{RegistroFrequencia}, e define a interface
\texttt{TipoAvaliacao}, com as operações \texttt{calcularNotaFinal()} e
\texttt{exigePresencaMinima()}, que abstrai os diferentes tipos de avaliação.
As regras de negócio são expostas pela interface
\texttt{RegistroAcademicoService}, implementada por
\texttt{RegistroAcademicoServiceImpl}, com operações como
\texttt{registrarNota()}, \texttt{consultarNotas()},
\texttt{registrarFrequencia()} e \texttt{consultarFrequencia()}. A
persistência é isolada pela interface \texttt{RegistroAcademicoRepository},
responsável por operações como \texttt{salvarNota()},
\texttt{buscarNotasPorAluno()}, \texttt{salvarFrequencia()} e
\texttt{buscarFrequenciaPorAluno()}. Essa separação entre serviço e
repositório mantém as regras de negócio independentes da tecnologia de
persistência, em linha com o objetivo de manutenibilidade (RNF-07).

O projeto de classes dos demais componentes será elaborado nas próximas
iterações, conforme cada um for detalhado.

\section{Princípios de projeto}

A equipe já tratava coesão e acoplamento como atributos de qualidade de
produto desde a etapa de especificação de requisitos, entendendo que, do
ponto de vista do desenvolvedor, a qualidade é medida pela implementação de
um código limpo, com alta coesão e baixo acoplamento, e associando esses
atributos diretamente à manutenibilidade (RNF-07). Já naquela etapa ficou
definido que a equipe monitoraria essas propriedades por meio de métricas
estáticas de produto, como o CBO (\emph{Coupling Between Object Classes}) e o
LCOM (\emph{Lack of Cohesion in Methods}), buscando manter CBO baixo e LCOM
próximo de zero.

Aplicado ao projeto de componentes do EduHub, esse princípio se traduz em
manter a separação de responsabilidades já estabelecida na divisão modular
do Capítulo 2: o componente de Controle de Acesso não deve conhecer as
regras de negócio de Gestão Acadêmica (por exemplo, cálculo de faltas ou
travamento de período letivo), e vice-versa. A relação entre eles deve se
limitar à consulta de permissão, verificando se o professor está autorizado a
lançar notas na turma. Da mesma forma, o componente de Relatórios e Auditoria
deve depender dos demais módulos apenas para leitura de dados, sem que módulos
como Gestão Acadêmica ou Cadastros precisem conhecer a lógica de geração de
relatórios. Essa baixa dependência mútua entre componentes é o que permite,
conforme RNF-07, adicionar novos recursos a um módulo sem impactar os
demais.

Além da coesão e do acoplamento, o projeto de classes do componente de
Autenticação e Acesso (Figura~\ref{fig:classes-autenticacao}) aplica três
princípios SOLID:
\begin{itemize}[nosep]
    \item \textbf{Responsabilidade Única (SRP)}: as entidades
    (\texttt{Usuario}, \texttt{Cargo}, \texttt{Permissao} e
    \texttt{HistoricoLogin}) representam apenas os dados e o estado, a interface
    \texttt{AutenticacaoService} concentra as regras de autenticação e de
    controle de acesso e a interface \texttt{Autenticacao\allowbreak Repository}
    cuida apenas da persistência.
    \item \textbf{Aberto/Fechado (OCP)}: a interface \texttt{Metodo2FA} deixa a
    autenticação de dois fatores aberta para extensão e fechada para
    modificação. Novos meios de verificação, como \texttt{TwoFactorEmail} e
    \texttt{TwoFactorApp}, são adicionados como classes concretas, e um
    eventual método por SMS entraria da mesma maneira, sem alterar o código
    existente.
    \item \textbf{Inversão de Dependência (DIP)}: a implementação
    \texttt{AutenticacaoServiceImpl} depende da abstração
    \texttt{Autenticacao\allowbreak Repository}, e não de uma tecnologia de
    persistência concreta, o que mantém as regras de controle de acesso
    independentes da infraestrutura.
\end{itemize}
Da mesma forma, o projeto de classes do componente de Registro Acadêmico
(Figura~\ref{fig:classes-registro}) aplica os mesmos três princípios:
\begin{itemize}[nosep]
    \item \textbf{Responsabilidade Única (SRP)}: cada classe tem um papel bem
    definido. As entidades (\texttt{Nota}, \texttt{Aluno}) representam apenas
    os dados e o estado, o serviço (\texttt{RegistroAcademicoService})
    concentra as regras de negócio e o repositório
    (\texttt{RegistroAcademicoRepository}) cuida apenas da persistência, o que
    facilita a manutenção e os testes isolados.
    \item \textbf{Aberto/Fechado (OCP)}: a interface \texttt{TipoAvaliacao}
    deixa o sistema aberto para extensão e fechado para modificação. Novos
    tipos de avaliação, como \texttt{AvaliacaoProva} e
    \texttt{AvaliacaoTrabalho}, são adicionados como classes concretas, sem
    alterar o código existente e sem recorrer a enumerações e a múltiplos
    condicionais.
    \item \textbf{Inversão de Dependência (DIP)}: a implementação
    \texttt{RegistroAcademicoServiceImpl} depende de abstrações, e não de
    detalhes técnicos. O \texttt{RegistroAcademicoRepository} oculta a
    tecnologia de persistência concreta, como MySQL, MongoDB ou uma API, o que
    mantém as regras de negócio independentes da infraestrutura.
\end{itemize}
Os princípios SOLID restantes, a Substituição de Liskov (LSP) e a Segregação
de Interface (ISP), e sua aplicação aos demais componentes serão detalhados
nas próximas iterações, conforme o projeto de classes de cada um for
elaborado.

\section{Padrões de projeto}

O projeto de classes apresentado na Seção~\ref{sec:detalhamento} e a divisão
em componentes do Capítulo 2 evidenciam padrões de projeto que se encaixam naturalmente nos problemas de
projeto do EduHub. Esta seção documenta os quatro padrões de encaixe mais
direto, indicando o problema que cada um resolve e um trecho de código
ilustrativo.

\textbf{Strategy (comportamental):} O componente de Registro Acadêmico
precisa calcular a nota final de diferentes tipos de avaliação, e o de
Autenticação e Acesso precisa validar diferentes métodos de segundo fator.
Tratar essas variações com uma cadeia de condicionais dentro do serviço
tornaria o código fechado à inclusão de novos tipos. O padrão Strategy
encapsula cada variação atrás de uma interface comum, deixando o serviço
independente da implementação concreta e permitindo acrescentar novos tipos
sem alterar o código existente, em conformidade com o princípio Aberto/Fechado.
Esse padrão já aparece no projeto de classes da Seção~\ref{sec:detalhamento},
nas interfaces \texttt{TipoAvaliacao} e \texttt{Metodo2FA}.

\begin{lstlisting}
interface TipoAvaliacao {
    double calcularNotaFinal(double nota);
    boolean exigePresencaMinima();
}
class AvaliacaoProva    implements TipoAvaliacao { /* ... */ }
class AvaliacaoTrabalho implements TipoAvaliacao { /* ... */ }

// o servico recebe a estrategia e nao conhece a implementacao concreta
double notaFinal = tipoAvaliacao.calcularNotaFinal(valor);
\end{lstlisting}

\textbf{Adapter (estrutural):} O EduHub envia notificações por meio de um
serviço de e-mail externo, cuja interface difere da esperada pelos componentes
internos. Como esse serviço é de terceiros e não pode ser modificado, o padrão
Adapter cria uma classe intermediária que implementa a interface interna
\texttt{Notificador} e traduz suas chamadas para a API externa. É esse o papel
do componente Adaptador de E-mail (RF-36).

\begin{lstlisting}
interface Notificador {
    void notificar(Usuario destino, String assunto, String corpo);
}

class AdaptadorEmail implements Notificador {
    private final ServicoEmailExterno externo; // API de terceiros, imutavel

    AdaptadorEmail(ServicoEmailExterno externo) { this.externo = externo; }

    public void notificar(Usuario destino, String assunto, String corpo) {
        // traduz a chamada interna para a API externa
        externo.send("no-reply@eduhub", destino.getEmail(), assunto, corpo);
    }
}
\end{lstlisting}

\textbf{Observer (comportamental):} Eventos como a publicação de um comunicado
(RF-35) ou o lançamento de uma nota exigem que vários destinos sejam
atualizados, como a caixa de entrada do aluno (RF-37) e o disparo de e-mail
(RF-36). Acoplar o produtor do evento a cada destino dificultaria a inclusão de
novos canais. O padrão Observer define uma relação um-para-muitos em que o
objeto observado notifica automaticamente todos os observadores registrados
quando seu estado muda.

\begin{lstlisting}
class EventoAcademico extends Subject {
    void publicarComunicado(Comunicado c) {
        this.comunicado = c;
        notifyObservers(); // avisa todos os canais registrados
    }
}
class CanalCaixaEntrada implements Observer {
    public void update(Subject s) { /* grava na caixa de entrada */ }
}
class CanalEmail implements Observer {
    public void update(Subject s) { /* dispara via AdaptadorEmail */ }
}
\end{lstlisting}

\textbf{Proxy de proteção (estrutural):} Regras como o acesso do professor
apenas às turmas às quais está vinculado (RF-21) e a restrição de alteração de
cargos a perfis administrativos (RF-07) exigem verificar a permissão antes de
executar a operação. Em vez de espalhar essa verificação pelas regras de
negócio, o padrão Proxy interpõe um objeto que implementa a mesma interface do
serviço real, verifica a permissão e apenas então delega a chamada.

\begin{lstlisting}
class RegistroAcademicoProxy implements RegistroAcademicoService {
    private final RegistroAcademicoService real;
    private final AutenticacaoService acesso;

    public Nota registrarNota(int alunoId, int disciplinaId,
                              double valor, TipoAvaliacao tipo) {
        if (!acesso.verificarPermissao(professorAtual(), "LANCAR_NOTA"))
            throw new AcessoNegado();
        return real.registrarNota(alunoId, disciplinaId, valor, tipo);
    }
}
\end{lstlisting}
